% titel: Blatt 9, A3: Benzol/Toluol-Dampfdruck bei 100°C
% Teil II: der durchgerechnete Weg.

\subsection*{Gegeben und gesucht}
Bei 100\,°C: $p^\circ_{\text{Benzol}}=1.801$ bar,
$p^\circ_{\text{Toluol}}=0.7310$ bar. Lösung siedet bei 1.000 bar.
Gesucht: Molenbrüche $x_B$, $x_T$.

\subsection*{Ansatz}
Ideale Lösung, Sieden bei $p_{\text{ges}}=1.000$ bar:
\[
  1.000 = x_B\cdot1.801 + x_T\cdot0.7310,\qquad x_B + x_T = 1 .
\]

\subsection*{Rechnung}
Einsetzen $x_T = 1-x_B$:
\[
  1.000 = x_B\cdot1.801 + (1-x_B)\cdot0.7310
  = 0.7310 + x_B\,(1.801-0.7310) .
\]
\[
  x_B = \frac{1.000-0.7310}{1.801-0.7310}
      = \frac{0.2690}{1.0700} \approx 0.251,
  \qquad x_T = 1-0.251 = 0.749.
\]

\begin{ergebnis}
$x_{\text{Benzol}} \approx 0.25$, $x_{\text{Toluol}} \approx 0.75$.
\end{ergebnis}

\begin{kontrolle}
Prüfung: $0.251\cdot1.801+0.749\cdot0.7310
=0.452+0.548=1.000$ bar. Probe stimmt. Da Benzol flüchtiger ist, hätte man
einen kleineren Anteil von Benzol erwartet als von Toluol --- passt.
\end{kontrolle}

\begin{merksatz}
Raoult: $p_{\text{ges}}=x_Ap_A^\circ+x_Bp_B^\circ$ mit $x_A+x_B=1$. Beim
Sieden gleichsetzen mit dem Aussendruck und nach dem Molenbruch auflösen.
\end{merksatz}
